
\documentclass[12pt,a4paper]{ctexart}

\usepackage[margin=2.3cm]{geometry}
\usepackage{amsmath}
\usepackage{array}
\usepackage{booktabs}
\usepackage{caption}
\usepackage{enumitem}
\usepackage{etoolbox}
\usepackage{fancyhdr}
\usepackage{float}
\usepackage{graphicx}
\usepackage{hyperref}
\usepackage{longtable}
\usepackage{tabularx}
\usepackage[table]{xcolor}
\usepackage{tikz}
\usetikzlibrary{arrows.meta,positioning,shapes.geometric,fit,calc}

\hypersetup{hidelinks}
\urlstyle{same}
\setlist{nosep}
\captionsetup{font=small}
\renewcommand{\arraystretch}{1.25}
\setlength{\parindent}{2em}
\setlength{\parskip}{0.3em}
\setlength{\tabcolsep}{4pt}
\emergencystretch=3em
\sloppy
\makeatletter
\g@addto@macro{\UrlBreaks}{\do\/\do\-\do\_\do\.\do\:}
\makeatother
\AtBeginEnvironment{longtable}{\small\setlength{\tabcolsep}{3pt}}
\AtBeginEnvironment{tabularx}{\small\setlength{\tabcolsep}{3pt}}
\AtBeginEnvironment{tabular}{\small\setlength{\tabcolsep}{3pt}}

\newcommand{\studentname}{许奕}
\newcommand{\studentid}{2351441}
\newcommand{\teamname}{许奕、王天一、马源胜、周由}
\newcommand{\coursename}{软件工程管理与经济}
\newcommand{\projectname}{基于大模型的建筑能源智能管理与运维优化系统}
\newcommand{\docversion}{V2.0}
\newcommand{\docdate}{2026年5月31日}
\newcommand{\code}[1]{{\footnotesize\nolinkurl{#1}}}
\newcolumntype{Y}{>{\raggedright\arraybackslash}X}
\newcolumntype{C}[1]{>{\centering\arraybackslash}p{#1}}
\newcolumntype{L}[1]{>{\raggedright\arraybackslash}p{#1}}

\tikzset{
  box/.style={rectangle, rounded corners, draw=cyan!55!black, fill=cyan!8, thick, minimum height=8mm, align=center, inner sep=5pt},
  process/.style={rectangle, rounded corners, draw=teal!55!black, fill=teal!8, thick, minimum height=8mm, align=center, inner sep=5pt},
  data/.style={cylinder, shape border rotate=90, aspect=0.25, draw=orange!70!black, fill=orange!10, thick, minimum height=12mm, align=center, inner sep=5pt},
  risk/.style={diamond, draw=red!60!black, fill=red!8, thick, aspect=2, align=center, inner sep=2pt},
  arrow/.style={-{Stealth[length=2.4mm]}, thick, draw=gray!70!black},
  note/.style={rectangle, draw=gray!60, fill=gray!6, rounded corners, align=left, inner sep=6pt}
}

\pagestyle{fancy}
\fancyhf{}
\fancyhead[L]{\coursename}
\fancyhead[R]{\reporttitle}
\fancyfoot[C]{\thepage}
\setlength{\headheight}{15pt}

\title{
    \vspace{-2em}
    \heiti \zihao{2} \reporttitle \\
    \vspace{1em}
    \zihao{4} 课程期末项目正式交付文档
}
\author{
    \songti \zihao{4}
    项目名称：\projectname \\
    项目组成员：\teamname \\
    负责人：\studentname \quad 学号：\studentid
}
\date{\docdate}

\newcommand{\reporttitle}{用户手册与部署说明}
\begin{document}

\maketitle
\thispagestyle{fancy}

\begin{abstract}
本文档为《\projectname》的用户手册与部署说明，面向系统使用者、评审人员和后续维护人员，说明如何获取代码、放置环境变量、安装依赖、启动后端、启动前端、启动 MCP Server、使用主要功能、处理常见问题和完成系统展示。
\end{abstract}

\tableofcontents
\newpage

\section{运行前准备}
\subsection{环境要求}
\begin{table}[H]
\centering
\caption{运行环境要求}
\begin{tabularx}{\textwidth}{p{4cm}Y}
\toprule
类别 & 要求 \\
\midrule
操作系统 & Windows 10/11，建议使用 PowerShell \\
Python & Python 3.11 或兼容版本 \\
Node.js & 支持 Vite 的 Node.js 版本 \\
浏览器 & Edge、Chrome 或其他现代浏览器 \\
网络 & 基础系统不强依赖外网；外部大模型需要网络和有效 API Key \\
\bottomrule
\end{tabularx}
\end{table}

\subsection{代码获取}
运行人员可从 GitHub 拉取代码：
\begin{verbatim}
git clone <GitHub 仓库地址>
cd building-energy-intelligence-platform
\end{verbatim}

\subsection{环境变量}
如果只演示基础系统，可复制模板：
\begin{verbatim}
Copy-Item .env.example .env
\end{verbatim}
如果需要启用外部大模型，应由项目负责人单独提供真实 \code{.env} 文件，并放在仓库根目录。真实 \code{.env} 不应放入 \code{backend/} 或 \code{frontend/}，不得提交到 GitHub，也不得放进最终 zip。

\section{启动后端}
\begin{verbatim}
python -m venv .venv
.venv\Scripts\Activate.ps1
pip install -r backend\requirements.txt
uvicorn app.main:app --reload --host 127.0.0.1 --port 8000 `
  --app-dir backend
\end{verbatim}
也可以直接运行脚本：
\begin{verbatim}
.\scripts\start-backend.ps1
\end{verbatim}
启动后可访问 \code{http://127.0.0.1:8000/docs} 查看接口文档，访问 \code{http://127.0.0.1:8000/api/v1/health} 检查健康状态。

\section{启动前端}
\begin{verbatim}
cd frontend
npm install
npm run dev
\end{verbatim}
也可以在项目根目录运行：
\begin{verbatim}
.\scripts\start-frontend.ps1
\end{verbatim}
前端默认访问地址为 \code{http://127.0.0.1:5173}。后端和前端通常需要两个 PowerShell 窗口分别运行。

\section{启动 MCP Server}
默认 stdio 模式：
\begin{verbatim}
.\scripts\start-mcp.ps1
\end{verbatim}
HTTP 模式：
\begin{verbatim}
.\scripts\start-mcp.ps1 -Transport streamable-http `
  -HostAddress 127.0.0.1 -Port 8765
\end{verbatim}
HTTP 模式地址为 \code{http://127.0.0.1:8765/mcp}。

\section{页面功能说明}
\subsection{总览}
总览页展示系统 KPI、时间范围、总体指标和三维楼层风险态势。用户可以拖拽三维视图观察不同角度，点击红色楼层进入统计分析。

\subsection{数据浏览}
数据浏览页用于查询原始和派生后的能耗记录。用户可选择建筑、楼层、开始时间、结束时间和显示数量。该页面用于全量查询，不受三维楼层点击联动影响。

\subsection{统计分析}
统计分析页是异常诊断工作台。用户可按建筑、楼层和时间范围锁定分析对象。若选中异常楼层，页面优先展示异常明细、异常原因、楼层设备、趋势和建议；若选中健康楼层，页面只展示健康提示。

\subsection{工单中心}
工单中心支持从异常记录创建工单。用户选择建筑、楼层、异常记录和负责人后，系统生成“处理中”工单。处理中工单显示在顶部，完成后变为绿色并靠后显示。工单状态保存到本地 JSON 文件，刷新页面不会丢失。

\subsection{决策报告}
决策报告页汇总当前范围内的能耗、异常、工单和优化建议，并提供简单收益模拟。该页面用于体现系统从异常发现到管理建议的综合价值。

\subsection{智能问答}
智能问答页支持本地知识库问答和可选外部模型增强。若外部模型不可用，系统会回退到本地知识库，保证演示不中断。

\section{推荐系统展示路径}
\begin{enumerate}
    \item 打开总览页，说明数据规模、建筑数量、时间范围和平均 COP。
    \item 拖拽三维楼宇，点击红色异常楼层，展示页面自动跳转统计分析。
    \item 在统计分析页查看异常明细，点击异常解释，说明触发规则和指标证据。
    \item 进入工单中心，从异常创建工单，分配负责人，完成工单并刷新验证持久化。
    \item 进入决策报告页，展示风险摘要、优化建议和收益模拟。
    \item 进入智能问答页，提问“当前有哪些异常风险”或“哪些设备优先维护”。
    \item 简要说明 MCP Server 可让 AI 客户端直接调用项目数据与分析工具。
\end{enumerate}

\section{常见问题}
\begin{longtable}{p{4cm}p{9.8cm}}
\caption{常见问题与处理}\\
\toprule
问题 & 处理方式 \\
\midrule
\endfirsthead
\toprule
问题 & 处理方式 \\
\midrule
\endhead
前端无法连接后端 & 确认后端运行在 8000 端口，Vite 代理配置未改动 \\
提示找不到数据文件 & 确认根目录存在 \code{data/samples/energy_records.csv}，不要从子目录启动错误脚本 \\
外部模型无回答 & 检查 \code{.env} 中 API Key、base URL、model 是否有效；也可关闭 LLM 使用本地回答 \\
工单刷新后丢失 & 检查 \code{data/runtime/work_orders.json} 是否可写 \\
npm install 慢 & 可更换 npm 镜像，但不要提交 \code{node_modules} \\
pytest 找不到模块 & 确认使用项目根目录运行检查脚本或正确设置 \code{PYTHONPATH} \\
\bottomrule
\end{longtable}

\clearpage
\section{补充用户操作手册}
\subsection{页面任务说明}
本节按用户实际操作任务组织，而不是按代码模块组织，便于第一次接触系统的同学或评审人员快速完成演示。
\begin{longtable}{L{2.4cm}L{3.2cm}L{5.8cm}L{3.8cm}}
\caption{用户任务、入口与操作结果}\\
\toprule
任务 & 入口 & 操作步骤 & 结果 \\
\midrule
\endfirsthead
\toprule
任务 & 入口 & 操作步骤 & 结果 \\
\midrule
\endhead
查看系统状态 & 总览页 & 打开首页，等待 KPI 和三维楼宇加载。 & 掌握整体能耗和风险。 \\
查看建筑风险 & 总览页 & 观察楼宇颜色和风险卡片。 & 发现异常楼层。 \\
旋转 3D 视图 & 总览页 & 按住鼠标拖拽楼宇区域。 & 从不同角度观察楼群。 \\
进入楼层分析 & 总览页 & 点击某个红色楼层。 & 跳转统计分析并锁定范围。 \\
查询原始记录 & 数据浏览 & 选择建筑、楼层、时间和数量后刷新。 & 查看符合条件的记录。 \\
导出记录 & 数据浏览 & 在筛选后点击导出。 & 下载 CSV 文件。 \\
查看异常明细 & 统计分析 & 选择建筑、楼层和时间范围。 & 优先看到异常记录。 \\
查看健康楼层 & 统计分析 & 选择无异常楼层。 & 只显示健康结论。 \\
查看趋势 & 统计分析 & 选择有异常的范围并查看趋势模块。 & 了解能耗变化。 \\
查看设备 & 统计分析 & 打开设备分析区域。 & 定位异常相关设备。 \\
解释异常 & 统计分析 & 点击异常解释或详情。 & 看到规则、证据和建议。 \\
创建工单 & 工单中心 & 选择建筑、楼层、异常和负责人。 & 生成处理中工单。 \\
完成工单 & 工单中心 & 点击完成按钮。 & 工单变绿并移动到后部。 \\
刷新验证 & 工单中心 & 刷新浏览器或重新进入页面。 & 工单状态仍保留。 \\
查看日报 & 决策报告 & 进入页面并查看风险摘要。 & 获得管理建议和收益模拟。 \\
使用问答 & AI 助手 & 选择模型或本地回答后输入问题。 & 获得能源运维回答。 \\
查看模型 & AI 助手 & 打开模型选择下拉框。 & 看到可用模型选项。 \\
查看 API & 后端 docs & 访问 \code{http://127.0.0.1:8000/docs}。 & 检查接口契约。 \\
启动 MCP & 命令行 & 执行 MCP 启动脚本。 & 为 AI 客户端提供工具。 \\
运行验证 & 项目根目录 & 执行 \code{scripts/check-project.ps1}。 & 确认测试和构建通过。 \\
\bottomrule
\end{longtable}

\subsection{常见问题排查}
\begin{longtable}{L{2.8cm}L{4.0cm}L{5.0cm}L{2.0cm}}
\caption{常见问题与处理方式}\\
\toprule
问题 & 可能原因 & 处理方式 & 优先级 \\
\midrule
\endfirsthead
\toprule
问题 & 可能原因 & 处理方式 & 优先级 \\
\midrule
\endhead
前端打不开 & Vite 服务未启动或端口不是 5173。 & 重新执行前端启动脚本并查看终端输出。 & 高 \\
接口报错 & 后端未启动或端口不是 8000。 & 启动后端并访问健康检查接口。 & 高 \\
数据为空 & 筛选条件过窄或时间范围无记录。 & 重置筛选或选择有数据的时间段。 & 中 \\
图表不显示 & 当前范围为健康楼层或无数据。 & 查看健康提示或扩大范围。 & 中 \\
工单不保存 & 运行目录无法写入。 & 检查 \code{data/runtime/} 权限。 & 高 \\
AI 没有外部回答 & 未配置真实 API Key。 & 使用本地知识库回答或补充 \code{.env}。 & 低 \\
模型请求失败 & 网络、额度或 Key 错误。 & 检查提供商配置并允许系统回退。 & 中 \\
MCP 启动失败 & 依赖或 transport 参数错误。 & 按数据接口文档检查命令。 & 中 \\
PDF 乱码 & 未使用 XeLaTeX 或字体环境异常。 & 使用脚本重新编译并确认 MiKTeX 可用。 & 中 \\
交付包太大 & 包含 \code{node_modules} 或 \code{dist}。 & 按交付说明删除依赖和构建产物。 & 高 \\
\bottomrule
\end{longtable}

\subsection{系统使用任务说明}
本节说明使用者完成安装、运行、主要功能操作和基础排查时需要覆盖的任务范围。
\begin{longtable}{L{1.8cm}L{3.2cm}L{6.4cm}L{4.0cm}}
\caption{系统使用任务说明}\\
\toprule
编号 & 任务 & 操作说明 & 完成结果 \\
\midrule
\endfirsthead
\toprule
编号 & 任务 & 操作说明 & 完成结果 \\
\midrule
\endhead
UT-01 & 环境准备 & 按照用户手册完成“环境准备”相关操作，遇到失败时先检查终端输出和常见问题表。 & 能够完成对应操作并获得预期页面或终端结果。 \\
UT-02 & 复制环境变量 & 按照用户手册完成“复制环境变量”相关操作，遇到失败时先检查终端输出和常见问题表。 & 能够完成对应操作并获得预期页面或终端结果。 \\
UT-03 & 安装后端依赖 & 按照用户手册完成“安装后端依赖”相关操作，遇到失败时先检查终端输出和常见问题表。 & 能够完成对应操作并获得预期页面或终端结果。 \\
UT-04 & 安装前端依赖 & 按照用户手册完成“安装前端依赖”相关操作，遇到失败时先检查终端输出和常见问题表。 & 能够完成对应操作并获得预期页面或终端结果。 \\
UT-05 & 启动后端 & 按照用户手册完成“启动后端”相关操作，遇到失败时先检查终端输出和常见问题表。 & 能够完成对应操作并获得预期页面或终端结果。 \\
UT-06 & 启动前端 & 按照用户手册完成“启动前端”相关操作，遇到失败时先检查终端输出和常见问题表。 & 能够完成对应操作并获得预期页面或终端结果。 \\
UT-07 & 打开总览 & 按照用户手册完成“打开总览”相关操作，遇到失败时先检查终端输出和常见问题表。 & 能够完成对应操作并获得预期页面或终端结果。 \\
UT-08 & 理解 KPI & 按照用户手册完成“理解 KPI”相关操作，遇到失败时先检查终端输出和常见问题表。 & 能够完成对应操作并获得预期页面或终端结果。 \\
UT-09 & 旋转楼宇 & 按照用户手册完成“旋转楼宇”相关操作，遇到失败时先检查终端输出和常见问题表。 & 能够完成对应操作并获得预期页面或终端结果。 \\
UT-10 & 点击楼层 & 按照用户手册完成“点击楼层”相关操作，遇到失败时先检查终端输出和常见问题表。 & 能够完成对应操作并获得预期页面或终端结果。 \\
UT-11 & 查询记录 & 按照用户手册完成“查询记录”相关操作，遇到失败时先检查终端输出和常见问题表。 & 能够完成对应操作并获得预期页面或终端结果。 \\
UT-12 & 筛选楼层 & 按照用户手册完成“筛选楼层”相关操作，遇到失败时先检查终端输出和常见问题表。 & 能够完成对应操作并获得预期页面或终端结果。 \\
UT-13 & 导出 CSV & 按照用户手册完成“导出 CSV”相关操作，遇到失败时先检查终端输出和常见问题表。 & 能够完成对应操作并获得预期页面或终端结果。 \\
UT-14 & 查看趋势 & 按照用户手册完成“查看趋势”相关操作，遇到失败时先检查终端输出和常见问题表。 & 能够完成对应操作并获得预期页面或终端结果。 \\
UT-15 & 查看异常 & 按照用户手册完成“查看异常”相关操作，遇到失败时先检查终端输出和常见问题表。 & 能够完成对应操作并获得预期页面或终端结果。 \\
UT-16 & 查看解释 & 按照用户手册完成“查看解释”相关操作，遇到失败时先检查终端输出和常见问题表。 & 能够完成对应操作并获得预期页面或终端结果。 \\
UT-17 & 查看设备 & 按照用户手册完成“查看设备”相关操作，遇到失败时先检查终端输出和常见问题表。 & 能够完成对应操作并获得预期页面或终端结果。 \\
UT-18 & 创建工单 & 按照用户手册完成“创建工单”相关操作，遇到失败时先检查终端输出和常见问题表。 & 能够完成对应操作并获得预期页面或终端结果。 \\
UT-19 & 完成工单 & 按照用户手册完成“完成工单”相关操作，遇到失败时先检查终端输出和常见问题表。 & 能够完成对应操作并获得预期页面或终端结果。 \\
UT-20 & 刷新验证 & 按照用户手册完成“刷新验证”相关操作，遇到失败时先检查终端输出和常见问题表。 & 能够完成对应操作并获得预期页面或终端结果。 \\
UT-21 & 查看日报 & 按照用户手册完成“查看日报”相关操作，遇到失败时先检查终端输出和常见问题表。 & 能够完成对应操作并获得预期页面或终端结果。 \\
UT-22 & 提出问题 & 按照用户手册完成“提出问题”相关操作，遇到失败时先检查终端输出和常见问题表。 & 能够完成对应操作并获得预期页面或终端结果。 \\
UT-23 & 选择模型 & 按照用户手册完成“选择模型”相关操作，遇到失败时先检查终端输出和常见问题表。 & 能够完成对应操作并获得预期页面或终端结果。 \\
UT-24 & 理解回退 & 按照用户手册完成“理解回退”相关操作，遇到失败时先检查终端输出和常见问题表。 & 能够完成对应操作并获得预期页面或终端结果。 \\
UT-25 & 启动 MCP & 按照用户手册完成“启动 MCP”相关操作，遇到失败时先检查终端输出和常见问题表。 & 能够完成对应操作并获得预期页面或终端结果。 \\
UT-26 & 查看 API 文档 & 按照用户手册完成“查看 API 文档”相关操作，遇到失败时先检查终端输出和常见问题表。 & 能够完成对应操作并获得预期页面或终端结果。 \\
UT-27 & 运行测试 & 按照用户手册完成“运行测试”相关操作，遇到失败时先检查终端输出和常见问题表。 & 能够完成对应操作并获得预期页面或终端结果。 \\
UT-28 & 编译文档 & 按照用户手册完成“编译文档”相关操作，遇到失败时先检查终端输出和常见问题表。 & 能够完成对应操作并获得预期页面或终端结果。 \\
UT-29 & 密钥安全 & 按照用户手册完成“密钥安全”相关操作，遇到失败时先检查终端输出和常见问题表。 & 能够完成对应操作并获得预期页面或终端结果。 \\
UT-30 & 整理交付包 & 按照用户手册完成“整理交付包”相关操作，遇到失败时先检查终端输出和常见问题表。 & 能够完成对应操作并获得预期页面或终端结果。 \\
UT-31 & 打开正式 PDF & 按照用户手册完成“打开正式 PDF”相关操作，遇到失败时先检查终端输出和常见问题表。 & 能够完成对应操作并获得预期页面或终端结果。 \\
UT-32 & 准备 PPT & 按照用户手册完成“准备 PPT”相关操作，遇到失败时先检查终端输出和常见问题表。 & 能够完成对应操作并获得预期页面或终端结果。 \\
UT-33 & 录制视频 & 按照用户手册完成“录制视频”相关操作，遇到失败时先检查终端输出和常见问题表。 & 能够完成对应操作并获得预期页面或终端结果。 \\
UT-34 & 系统展示 & 按照用户手册完成“系统展示”相关操作，遇到失败时先检查终端输出和常见问题表。 & 能够完成对应操作并获得预期页面或终端结果。 \\
UT-35 & 故障排查 & 按照用户手册完成“故障排查”相关操作，遇到失败时先检查终端输出和常见问题表。 & 能够完成对应操作并获得预期页面或终端结果。 \\
\bottomrule
\end{longtable}

\end{document}
